Our proposed modularized solution for learning-based navigation builds upon a \ac{vae}~\cite{doersch2016tutorial,kingma2013auto} module that can be independently trained using both real and simulated data to produce a compressed latent representation of a raw depth image that is then exploited for collision prediction. Semantic labels can also be added to the images --when available-- to ensure reconstruction of hard-to-perceive thin obstacles (e.g., skinny rods and wires~\cite{Kulkarni2020}). This representation is able to reconstruct complex images without missing out on narrow cross-section (thin) obstacles, while being robust to sensor noise that arises from the sensor's inability to reconstruct depth for textureless or reflective regions and stereo shadows. This latent representation is then used to train a $3$D \ac{cpn} in simulation to predict collision scores for candidate action sequences based on a library of motion primitives. Building upon our prior work called ORACLE~\cite{ORACLE_ICRA2022}, the method does not require the full robot state or an environment map but only real-time depth frames and access to an estimate of the partial robot state involving the system's attitude, as well as angular and linear velocities. Simultaneously, it takes into account the uncertainty associated with this partial state when predicting collision scores. Given the above and a goal vector provided by a high-level planner~\cite{GBPLANNER_JFR_2020,achtelik2014motion} or by an operator, the optimal collision-free action sequence is selected and executed by the robot in a receding horizon manner.


\subsection{Semantically-enhanced Variational Autoencoder}\label{sec:vae}
\ac{vae}s are powerful tools that can allow encoding high-dimensional input data in a compressed representation. In this work, we show that a complex depth image with a typical resolution (of $270\times 480$ pixels) can be sufficiently represented by a highly compressed latent representation (here $128$ variables) while preserving features from hard-to-perceive thin obstacles. This latent representation can subsequently be used for predicting collisions given a set of candidate action sequences. We develop an approach to combine data from both simulated and real sensor observations, allowing our method to be a) robust to sensor noise, and at the same time b) able to utilize additional information such as instance semantic labels --primarily from simulators-- thus being capable of reconstructing thin objects. In contrast, non-modularized methods that train a collision predictor end-to-end such as~\cite{ORACLE_ICRA2022} may not allow to explicitly promote a focus on such semantically-driven thin obstacles or conveniently utilize real sensor data and their imperfections.

%%%%%%%%%%%%%%%%%%%%%%%%%%%%%%%%%%%%%%%%%%%%%%%%%%%%%%%%%%%%%%%%%
\begin{figure*}[ht]
\includegraphics[width=\textwidth]{figures/VAE_arch_only_seVAE_v2_corrected.png}
\vspace{-4ex}
\caption{Proposed architecture for the \ac{sevae} for collision prediction. The convolutional and deconvolutional layers' hyperparameters are represented in the format ($\mathfrak{a}$ $\times$ $\mathfrak{b}$ (de)conv$\mathfrak{, c, /d}$), where $\mathfrak{a}$ $\times$ $\mathfrak{b}$ refers to the kernel size, $\mathfrak{c}$ refers to the number of channels, and $\mathfrak{d}$ refers to the stride length. The dense layers only have the layer size mentioned alongside. A latent space $\mathbf{z}$ of $128$ dimensions is sampled from the distribution with mean $\bm{\mu}$ and std. dev. $\bm{\sigma}$ and used to reconstruct the image by the decoder.}
\vspace{-4ex}
\label{fig:VAE_network_architecture}
\end{figure*}
%%%%%%%%%%%%%%%%%%%%%%%%%%%%%%%%%%%%%%%%%%%%%%%%%%%%%%%%%%%%%%%%%


Representing a compressed latent space starting from high-dimensional depth data using hand-tuned parameters is difficult, given that the distribution generating the dataset of depth images is practically intractable to compute. Thus, we utilize a \ac{vae} to learn a latent representation that can be used for image reconstruction. We consider a dataset $\mathbf{X} = \{\mathbf{x}^{(i)}, \mathbf{x}_{val}^{(i)}, \mathbf{x}_{seg}^{(i)}\}^N_{i=1}$ consisting of samples of discrete images $\mathbf{x}^{(i)}$ of dimensions $H\times W$ (here $270\times480$ pixels), the mask of valid pixels from the input data $\mathbf{x}_{val}^{(i)}$, and the instance segmentation mask for each thin obstacle instance $\mathbf{x}_{seg}^{(i)}$ (when available). The pixels with defined depth information from the sensor are referred to as valid pixels and the pixels with missing depth values are referred to as invalid pixels. We assume that this dataset is generated by some random process involving an unobserved random variable $\mathbf{z}$ with $J$ dimensions (here $J=128$). We model this variable $\mathbf{z}$ as the latent code for our representation. We assume that a depth image $\mathbf{x}$ is generated from $\mathbf{z}$ by some generative model. A probabilistic decoder $p_{\bm{\theta}}(\mathbf{x}|\mathbf{z})$ produces a distribution over the possible values of $\mathbf{x}$. Since the true value of the posterior $p_{\bm{\theta}}(\mathbf{z}|\mathbf{x})$ is unknown, we assume an approximate posterior $q_{\bm{\phi}}(\mathbf{z}| \mathbf{x})$, to be a multivariate Gaussian with a diagonal covariance as:

\small
\begin{equation}
    \log q_{\bm{\phi}}(\mathbf{z}| \mathbf{x}^{(i)}) = \log \mathcal{N}(\mathbf{z}; \bm{\mu}^{(i)}, {\bm{\sigma}}^{2(i)}\mathbf{I}),
\end{equation}
\normalsize
where the mean and the standard deviation of the approximate posterior, $\bm{\mu}$ and $\bm{\sigma}$, are outputs of the encoding neural network for the image $\mathbf{x}$ with parameters $\bm{\phi}$. In this realization, we use the reparameterization trick to obtain the values of $\bm{\mu}^{(i)}$ and $\log({\bm{\sigma}}^{2(i)})$ from the network~\cite{doersch2016tutorial}. We then sample $\mathbf{z}^{(i)} \sim q_{\bm{\phi}}(\mathbf{z} | \mathbf{x}^{(i)})$ using $\mathbf{z}^{(i)} = \bm{\mu}^{(i)} + \bm{\sigma}^{(i)}\odot\bm{\epsilon}$, ($\odot$ is the element-wise multiplication operator), with an auxiliary noise variable $\bm{\epsilon} \sim \mathcal{N}(\mathbf{0},\mathbf{I})$.  A suitable encoder for this architecture is designed by modifying~\cite{loquercio2018dronet} to reduce the number of parameters in the neural network, while enhancing connectivity in the fully-connected layer. The decoder consists of deconvolutional layers stacked with activation functions without skip connections. The network architecture used in this work is depicted in Figure~\ref{fig:VAE_network_architecture}, where its fine tuning was guided by the observation that increasing the size of the fully connected hidden layers tended to enhance the quality of reconstruction from the network. The loss term is defined as:

\small
\begin{equation}
    \mathcal{L} = \mathcal{L}_{recon} + \beta_{norm}\mathcal{L}_{KL},
\end{equation}
\normalsize
where $\mathcal{L}_{KL}$ denotes the KL-divergence loss defined as:

\small
\begin{equation}
    \mathcal{L}_{KL}(\bm{\mu}, \bm{\sigma}) = \frac{1}{2}\sum_{j=1}^J \left (1 + \log((\bm{\sigma}_j)^2) - (\bm{\mu}_j)^2 - (\bm{\sigma}_j)^2 \right).
\end{equation}
\normalsize
The constant $\beta_{norm} = \frac{\beta\cdot J}{H\cdot W}$ weighing the KL-divergence loss is normalized accounting for the latent space and image dimensions. The value of $\beta$ (here $\beta = 1$) allows the mixing of the contributions from the KL-divergence loss and the reconstruction loss and is a tunable hyperparameter~\cite{higgins2017betavae}. The depth cameras used to collect the datasets consist of a stereo pair and are unable to reconstruct depth from textureless features and miss the depth data around the edges of features resulting in a ``stereo shadow'' region. We design a semantic-weighted reconstruction loss function $\mathcal{L}_{recon}$ for training the neural network taking into consideration these errors from such sensors by ignoring the contribution of the invalid pixels in the loss function. The semantic-weighted reconstruction loss is defined as:
\vspace{-1ex}
\begin{multline}\label{eq:reconstruction_loss}
    \mathcal{L}_{recon}(\mathbf{x}, \mathbf{x}_{recon}, \mathbf{x}_{val}, \mathbf{x}_{seg}) =\\ \sum_{h=1}^{H}\sum_{w=1}^{W}((\mathbf{x}-\mathbf{x}_{recon})^2\odot{\mathbf{x}_{val}}\odot{\lambda(\mathbf{x}_{seg})})_{(h,w)},
\end{multline}
where $\mathbf{x}_{recon}$ is the reconstructed image and $\{\mathbf{x}, \mathbf{x}_{val}, \mathbf{x}_{seg}\}$ represents one sample from the dataset $\mathbf{X}$. The use of $\mathbf{x}_{val}$ during \ac{sevae} training eliminates the contribution of the invalid input pixels from the loss function, allowing the network to represent distributions consisting of information in the valid pixels. The function $\lambda(\mathbf{x}_{seg})$ creates a pixel-wise weight mask to give higher weight to the pixels that correspond to disproportionately thin obstacles using the semantic label information $\mathbf{x}_{seg}$ (if available). The weight $\nu_\ell$ of a pixel $\ell$ belonging to a semantic instance $S_k \in \mathbb{S}$, where $\mathbb{S}$ is the set of all instances of semantics in $\mathbf{x}_{seg}$, depends on the pixel count $p_k$ of instance $S_k$. The weight of the pixels not belonging to a semantic instance is set to 1. Formally, the weight $\nu_\ell$ is defined as:

\small
\begin{equation}
\nu_\ell = 
    \begin{cases}
    \max(W_{const}/p_k, \nu_{\min}),~ \ell \in S_k~\textrm{and}~p_k > p_{\min}\\
    1,~\textrm{otherwise}
    \end{cases},
\end{equation}
\normalsize
where the term $W_{const}$ (here $W_{const}=6000$) acts as a multiplicative constant to weigh the inverse count of pixels per semantic, while $\nu_{\min}$ (here $\nu_{\min}=15$) limits the minimum per-pixel weight. This weighing term is applied to magnify the contribution of small-sized semantics while allowing larger-sized semantics to be proportionally weighed based on the number of pixels. We ignore semantics smaller than $p_{\min}$ (here $p_{\min}=40$ pixels) to prevent the \ac{sevae} from trying to reconstruct extremely small regions, that may, in the real-world correspond to sensor noise and data imperfections. Given the loss $\mathcal{L}$ and the semantic information, the \ac{sevae} is trained to encode a depth image into the low-dimensional latent representation $\boldsymbol{\mu}$ to be used for collision prediction.


\subsection{Collision Prediction Network and Action Planning}\label{sec:collision_prediction_network}

The latent space from the \ac{sevae} is then exploited by the \ac{cpn} to facilitate safe navigation building upon our previous work ORACLE~\cite{ORACLE_ICRA2022} and extending it to $3$D navigation, while retaining the benefits of the modularized \ac{sevae}-based latent space including its ability to assimilate real sensor data. To this end, the CNN part of the original \ac{cpn} in ORACLE is replaced by the \ac{sevae} which takes the current depth image $\mathbf{x}_t$ and outputs the compressed latent vector $\boldsymbol{\mu}_t$, that is then fed to the \ac{cpn} as illustrated in Figure~\ref{fig:CPN_network_architecture}.

Specifically, let $\Bs,\Vs$ be the body frame and vehicle frame of the robot respectively, and $\mathbf{s}_t = [{\mathbf{v}_t}^T, \omega_t, \vartheta_t, \varphi_t]^T $ the estimated partial state of the robot at time $t$ consisting of a) the $3$D velocity in $\Vs$ ($\mathbf{v}_t=[v_{t,x}, v_{t,y}, v_{t,z}]^T \in \mathbb{R}^{3 \times 1}$), b) the angular velocity around the $z$-axis of the body-frame $\Bs$ ($\omega_t$), as well as c) the roll ($\vartheta_t$) and pitch angles ($\varphi_t$). Let $\boldsymbol{\Sigma}_t$ denote the covariance of the estimated robot’s partial state, $\mathbf{n}^g_t$ the $3$D unit goal vector, expressed in $\Vs$, that the robot has to follow, $\psi_t$ the current yaw angle of the robot, and $\mathbf{a}_{t:t+T}=[\mathbf{a}_t, \mathbf{a}_{t+1},\hdots,\mathbf{a}_{t+T-1}]$ an action sequence having length $T$ where the action at time step $t+i~(i=0,\hdots,T-1)$ includes a) the reference $3$D velocity expressed in the vehicle frame $\mathbf{v}^r_{t+i}$ and
b) the steering angle ($\delta^r_{t+i}$) from the current yaw angle of the robot ($\psi_{t}$), such that $\mathbf{a}_{t+i}=[{\mathbf{v}^r_{t+i}}^T, \delta^r_{t+i}]^T$.
The method finds an optimized collision-free sequence of actions $\mathbf{a}_{t:t+T}$ given $(\mathbf{x}_t,\mathbf{s}_t,\boldsymbol{\Sigma}_t)$ (specifically using the latent vector $\bm{\mu}_t$ derived from $\mathbf{x}_t$ using the \ac{sevae}'s encoder) enabling the robot to safely navigate to the goal vector $\mathbf{n}^g_t$. 

The \ac{cpn} processes $\bm{\mu}_t$, $\mathbf{s}_t$, and a Motion Primitives-based Library (MPL) of $N_{MP}$ sequences of future velocity and steering angle references $\mathbf{a}_{t:t+T}$ to predict the collision scores $\hat{\mathbf{c}}^{col}_{t+1:t+T+1}=[\hat{c}^{col}_{t+1}, \hat{c}^{col}_{t+2},\hdots,\hat{c}^{col}_{t+T}]$ of the robot at each time step from $t+1$ to $t+T$ in the future for each action sequence. 
To account for the uncertainty in the robot's partial state, we first calculate $N_{\Sigma} = 2\gamma + 1$ sigma points ($\mathbf{m}_1=\mathbf{s}_t$, $\hdots$, $\mathbf{m}_{N_{\Sigma}}$), where $\gamma$ is the dimension of the robot's partial state (here $\gamma=6$), based on the mean $\mathbf{s}_t$ and covariance $\bm{\Sigma}_t$ using the Unscented Transform (UT) \cite{Julier1997unscented}. Using the sigma points, the uncertainty-aware collision score $\hat{c}^{uac}$ is calculated as presented in~\cite{ORACLE_ICRA2022}. A set of safe action sequences is derived by thresholding $\hat{c}^{uac}$ and the action sequence that leads to the end velocity of the robot best aligned with $\mathbf{n}^g_t$ is chosen. The first action in the sequence is executed, while the process is repeated in a receding horizon fashion. Notably, in ORACLE, the \ac{cpn} is trained end-to-end entirely in simulation and requires a computationally-expensive pre-processing step in the depth image input to close the sim-to-real gap. In this work, the \ac{cpn} is trained with the frozen weights of the \ac{vae} in simulation environments which have large obstacles with randomized shapes and thin obstacles with diameters less than $5~\textrm{cm}$, as illustrated in Figure~\ref{fig:data_collection_diagram}.3.

%%%%%%%%%%%%%%%%%%%%%%%%%%%%%%%%%%%%%%%%%%%%%%%%%%%%%%%%%%%%%%%%%
\begin{figure}[h]
\includegraphics[width=\columnwidth]{figures/CPN_architecture_v5.png}
\vspace{-4ex}
\caption{Proposed architecture for the modularized approach involving the \ac{sevae} and the Collision Prediction Network (CPN). The \ac{sevae} encodes the input depth image $\mathbf{x}_t$ into the latent representation $\bm{\mu}_t$ which is used by the \ac{cpn} to predict the collision scores $\hat{\mathbf{c}}^{col}_{t+1:t+T+1}$ for each action sequence $\mathbf{a}_{t:t+T}$ in the motion primitives library. Furthermore, the method utilizes the $N_{\Sigma}$ sigma points calculated based on $\mathbf{s}_t$ and $\bm{\Sigma}_t$ to calculate the robot's partial state uncertainty-aware collision score $\hat{c}^{uac}$.}
\vspace{-2ex}
\label{fig:CPN_network_architecture}
\end{figure}
%%%%%%%%%%%%%%%%%%%%%%%%%%%%%%%%%%%%%%%%%%%%%%%%%%%%%%%%%%%%%%%%%
